\section*{Chapter \ref{ch:b1:poly} --- Polynomials}

\begin{solution}{exo:b1:poly:1}
$X^5 - 1 = (X^2 + X + 1)(X^3 - X^2 + 1) + (-X - 2)$. Steps: subtract
$X^3 B$, then $-X^2 B$, then $B$; the remainder $-X - 2$ has degree
$1 < 2$. \emph{Check at $X = 1$:} $\;0 = 3 \times 1 + (-3)$.

$2X^4 + X^3 - X + 3 = (X^2 - 2)(2X^2 + X + 4) + (X + 11)$. \emph{Check
at $X = 0$:} $\;3 = (-2)(4) + 11$.
\end{solution}

\begin{solution}{exo:b1:poly:2}
$X^2 + X + 1 = (X - j)(X - j^2)$ with $j = \eu^{2\iu\pi/3}$, $j^3 =
1$. It \omterm{def:b1:arith:divides}{divides} $Q_n = X^{2n} + X^n + 1$ iff $j$ and $j^2$ are roots of
$Q_n$; since $Q_n$ has real coefficients, $Q_n(j^2) =
\conj{Q_n(j)}$, so the condition is just $Q_n(j) = 0$. Now $Q_n(j) =
j^{2n} + j^n + 1$ depends on $n$ mod $3$:
\begin{itemize}
  \item $n \equiv 0$: $Q_n(j) = 1 + 1 + 1 = 3 \neq 0$;
  \item $n \equiv 1$: $Q_n(j) = j^2 + j + 1 = 0$;
  \item $n \equiv 2$: $Q_n(j) = j^4 + j^2 + 1 = j + j^2 + 1 = 0$.
\end{itemize}
So $X^2 + X + 1 \mid X^{2n} + X^n + 1$ exactly when $3 \nmid n$.
\end{solution}

\begin{solution}{exo:b1:poly:3}
By \cref{prop:b1:poly:multiplicity}, $(X-1)^2 \mid P$ iff $P(1) =
P'(1) = 0$:
\[
P(1) = 2 + a + b = 0, \qquad P'(1) = 4 + 3a + 2b = 0 .
\]
Solving: $b = -a - 2$ and $4 + 3a - 2a - 4 = a = 0$, so $a = 0$, $b =
-2$: $P = X^4 - 2X^2 + 1 = (X^2 - 1)^2 = (X-1)^2 (X+1)^2$, which is
the real factorization.
\end{solution}

\begin{solution}{exo:b1:poly:4}
$X^3 - 1 = (X - 1)(X - j)(X - j^2)$ over $\C$ ($j = \eu^{2\iu\pi/3}$),
and $(X - 1)(X^2 + X + 1)$ over $\R$.

$X^4 + X^2 + 1 = (X^2 + X + 1)(X^2 - X + 1)$ over $\R$ (multiply out,
or note $X^4 + X^2 + 1 = (X^2+1)^2 - X^2$); over $\C$, each quadratic
splits: roots $j, j^2$ and $-j, -j^2$, i.e.\ $\eu^{\pm 2\iu\pi/3},
\eu^{\pm\iu\pi/3}$.

$X^6 - 1 = \prod_{k=0}^{5} (X - \eu^{\iu k\pi/3})$ over $\C$, and over
$\R$:
\[
X^6 - 1 = (X-1)(X+1)(X^2 + X + 1)(X^2 - X + 1),
\]
grouping the conjugate pairs $\eu^{\pm 2\iu\pi/3}$ and $\eu^{\pm
\iu\pi/3}$.
\end{solution}

\begin{solution}{exo:b1:poly:5}
\begin{enumerate}
  \item Let $p/q$ (lowest terms) be a root of the \omterm{def:b1:poly:def}{monic} integer
        \omterm{def:b1:poly:def}{polynomial} $X^3 + \dots + c_0$: clearing denominators in
        $P(p/q) = 0$ gives $p^3 = -q\,(\text{integer})$, so $q \mid
        p^3$; coprimality forces $q = \pm 1$: the root is an integer
        $p$, and $p \mid c_0$ (isolate $c_0$). Here the candidates
        divide $6$: testing, $P(1) = 0$, $P(2) = 0$, $P(3) = 0$. So
        $P = (X-1)(X-2)(X-3)$.
  \item Vieta: $s_1 = 6$, $s_2 = 11$, $s_3 = 6$. Sum of squares:
        $s_1^2 - 2s_2 = 36 - 22 = 14 = 1 + 4 + 9$, as expected. Sum of inverses:
        $\frac{s_2}{s_3} = \frac{11}{6} = 1 + \frac12 + \frac13$, as expected.
\end{enumerate}
\end{solution}

\begin{solution}{exo:b1:poly:6}
First division step of the \omterm{met:b1:arith:euclid}{Euclidean algorithm}:
\[
(X + 1)(X^3 - X^2 + X - 1)
= X^4 - X^3 + X^2 - X + X^3 - X^2 + X - 1 = X^4 - 1 ,
\]
so the division of $X^4 - 1$ by $X^3 - X^2 + X - 1$ is exact
(quotient $X + 1$, remainder $0$), and the algorithm stops at once:
\[
\gcd(X^4 - 1,\; X^3 - X^2 + X - 1) = X^3 - X^2 + X - 1
\]
(already \omterm{def:b1:poly:def}{monic}). The Bézout relation is the trivial one: $\gcd = 0
\cdot (X^4 - 1) + 1 \cdot (X^3 - X^2 + X - 1)$. Consistency check by
factoring: $X^3 - X^2 + X - 1 = (X - 1)(X^2 + 1)$, which is indeed
the product of the common irreducible factors of $X^4 - 1 =
(X-1)(X+1)(X^2+1)$.
\end{solution}

\begin{solution}{exo:b1:poly:7}
Since $P \geq 0$ on $\R$, its real roots have even \omterm{def:b1:poly:derivative}{multiplicity} (at a
root of odd \omterm{def:b1:poly:derivative}{multiplicity}, $P$ changes sign). Using
\cref{cor:b1:poly:factorization} and pairing, write
\[
P = c \prod_i (X - a_i)^{2k_i} \prod_j \bigl((X - z_j)(X - \conj
z_j)\bigr)^{n_j},
\]
with $c > 0$ (behavior at $+\infty$). Let
\[
S = \sqrt c\, \prod_i (X - a_i)^{k_i} \prod_j (X - z_j)^{n_j}
\in \C[X],
\]
so that $P = S\,\conj S$ where $\conj S$ has the conjugated
coefficients. Split $S = A + \iu B$ with $A, B \in \R[X]$: then
\[
P = (A + \iu B)(A - \iu B) = A^2 + B^2 .
\]
\end{solution}

\begin{solution}{exo:b1:poly:8}
Lagrange (\cref{thm:b1:poly:lagrange}) with nodes $0, 1, 2$:
\[
P = 1\cdot\frac{(X-1)(X-2)}{(0-1)(0-2)} + 3\cdot\frac{X(X-2)}{1\cdot(1-2)}
+ 2\cdot\frac{X(X-1)}{2\cdot 1}
= \frac{(X-1)(X-2)}{2} - 3X(X-2) + X(X-1).
\]
Expanding: $\frac{X^2 - 3X + 2}{2} - 3X^2 + 6X + X^2 - X =
-\frac{3}{2}X^2 + \frac{7}{2}X + 1$.

System: $P = aX^2 + bX + c$ with $c = 1$; $a + b + 1 = 3$; $4a + 2b +
1 = 2$. Subtracting twice the second from the third: $2a - 1 = -4$,
so $a = -\frac32$, $b = \frac72$. Same \omterm{def:b1:poly:def}{polynomial}:
$P = -\frac32 X^2 + \frac72 X + 1$. (Check $P(2) = -6 + 7 + 1 = 2$.)
\end{solution}

\begin{solution}{exo:b1:poly:9}
$P = X^{2n+1} - 1$: $P' = (2n+1)X^{2n} \geq 0$, so the \omterm{def:b1:poly:def}{polynomial}
function is increasing (strictly except at $0$), with limits
$\mp\infty$: it vanishes exactly once on $\R$ (at $x = 1$).

Let $E_n = \sum_{k=0}^{n} \frac{X^k}{k!}$. Then $E_n' = E_{n-1} = E_n
- \frac{X^n}{n!}$. A multiple root $a$ would satisfy $E_n(a) =
E_n'(a) = 0$ (\cref{prop:b1:poly:multiplicity}), hence
$\frac{a^n}{n!} = E_n(a) - E_n'(a) = 0$, so $a = 0$; but $E_n(0) = 1
\neq 0$. No multiple root.
\end{solution}

\begin{solution}{exo:b1:poly:10}
\begin{enumerate}
  \item Induction (two base cases hold). Using $\cos(n+1)\theta +
        \cos(n-1)\theta = 2\cos\theta\cos n\theta$:
        \[
        T_{n+1}(\cos\theta) = 2\cos\theta \cos n\theta -
        \cos(n-1)\theta = \cos(n+1)\theta .
        \]
  \item $T_n(\cos\theta) = 0$ iff $n\theta \equiv \frac\pi2 \pmod
        \pi$: the numbers
        \[
        x_k = \cos\Bigl(\frac{(2k+1)\pi}{2n}\Bigr),
        \qquad k = 0, 1, \dots, n-1,
        \]
        are $n$ distinct points of $\intoo{-1}{1}$ (the angles lie in
        $\intoo{0}{\pi}$ where $\cos$ is \omterm{def:b1:logic:inj}{injective}), all roots of
        $T_n$; since $\deg T_n = n$ (from the recurrence, with
        leading coefficient $2^{n-1}$ for $n \geq 1$, by induction),
        these are \emph{all} the roots, each simple.
  \item For $x = \cos\theta \in \intcc{-1}{1}$: $\abs{T_n(x)} =
        \abs{\cos n\theta} \leq 1$, with equality iff $n\theta \equiv
        0 \pmod\pi$, i.e.\ at the $n+1$ points $y_k =
        \cos\frac{k\pi}{n}$, $k = 0, \dots, n$, where $T_n(y_k) =
        (-1)^k$. (This equioscillation makes $2^{1-n}T_n$ the \omterm{def:b1:poly:def}{monic}
        degree-$n$ \omterm{def:b1:poly:def}{polynomial} of smallest sup-norm on
        $\intcc{-1}{1}$ --- proved in this chapter's weekend
        problem.)
\end{enumerate}
\end{solution}

\begin{solution}{exo:b1:poly:11}
Write $P = c\prod_i (X - a_i)^{m_i}$. The product rule (extended to
several factors) gives
\[
P' = c\sum_{i} m_i (X - a_i)^{m_i - 1} \prod_{k \neq i} (X -
a_k)^{m_k},
\]
and dividing by $P$: $\frac{P'}{P} = \sum_i \frac{m_i}{X - a_i}$
(as rational functions, i.e.\ away from the roots).

Let $w$ be a root of $P'$. If $w$ is one of the $a_i$, it lies in the
convex hull trivially. Otherwise, evaluating at $w$:
\[
0 = \sum_i \frac{m_i}{w - a_i}
= \sum_i m_i\, \frac{\conj w - \conj a_i}{\abs{w - a_i}^2} .
\]
Conjugating: $\sum_i \lambda_i (w - a_i) = 0$ where $\lambda_i =
\frac{m_i}{\abs{w - a_i}^2} > 0$. Hence
\[
w = \frac{\sum_i \lambda_i a_i}{\sum_i \lambda_i} :
\]
a convex combination (positive weights summing to $1$ after
normalization) of the roots $a_i$. So every root of $P'$ lies in the
convex hull of the roots of $P$.
\end{solution}

\begin{solution}{exo:b1:poly:12}
Sum the evaluations of $(1 + X)^n$ at the three cube \omterm{def:b1:complex:unity}{roots of
unity}:
\[
2^n + (1 + j)^n + (1 + j^2)^n
= \sum_{k=0}^n \binom nk\,\bigl(1 + j^k + j^{2k}\bigr)
= 3\sum_{k\,:\,3\mid k}\binom nk ,
\]
since $1 + j^k + j^{2k}$ is a geometric sum equal to $3$ when $3
\mid k$ and to $\frac{j^{3k} - 1}{j^k - 1} = 0$ otherwise. Now $1 +
j = \frac12 + \iu\frac{\sqrt3}2 = \eu^{\iu\pi/3}$ and $1 + j^2 =
\conj{1 + j} = \eu^{-\iu\pi/3}$, so $(1+j)^n + (1+j^2)^n =
2\cos\frac{n\pi}3$ and
\[
\sum_{k\geq0}\binom n{3k} = \frac{2^n + 2\cos\frac{n\pi}3}{3} .
\]
Checks: $n = 3$: $\frac{8 + 2\cos\pi}3 = 2 = \binom30 + \binom33$;
$n = 6$: $\frac{64 + 2}3 = 22 = 1 + 20 + 1$.
\end{solution}

\begin{solution}{pb:b1:poly:1}
\textbf{1.} $T_2 = 2X^2 - 1$; $T_3 = 2X(2X^2 - 1) - X = 4X^3 -
3X$; $T_4 = 2X\,T_3 - T_2 = 8X^4 - 8X^2 + 1$; $T_5 = 2X\,T_4 - T_3
= 16X^5 - 20X^3 + 5X$. The identity $T_3(\cos\theta) = \cos3\theta$
is exactly $\cos3\theta = 4\cos^3\theta - 3\cos\theta$ from
\cref{ex:b1:complex:cos3}.

\textbf{2.} True for $n = 1, 2$. If $T_{n-1}$, $T_n$ have degrees
$n-1$, $n$ and leading coefficients $2^{n-2}$, $2^{n-1}$, then
$2X\,T_n$ has degree $n+1$ and leading coefficient $2^n$, while
$T_{n-1}$ has lower degree: $T_{n+1}$ has degree $n + 1$, leading
coefficient $2^n$. Parity: if $T_{n-1}$ has the parity of $n - 1$
and $T_n$ that of $n$, then $2X\,T_n$ and $T_{n-1}$ both have the
parity of $n + 1$, hence so does $T_{n+1}$.

\textbf{3.} If $P(\cos\theta) = \cos n\theta$ for all $\theta$,
then $P$ and $T_n$ agree at every point of $\intcc{-1}1$ --- an
infinite \omterm{def:b1:logic:sets}{set} --- so $P - T_n$ has infinitely many roots and is the
zero \omterm{def:b1:poly:def}{polynomial} (\cref{cor:b1:poly:nroots}).

\textbf{4.} For $x = \cos\theta$: $T_m(T_n(\cos\theta)) =
T_m(\cos n\theta) = \cos(mn\theta) = T_{mn}(\cos\theta)$, and
$2T_mT_n(\cos\theta) = 2\cos m\theta\cos n\theta = \cos(m+n)\theta
+ \cos\abs{m - n}\theta$. Both identities hold on $\intcc{-1}1$,
hence as \omterm{def:b1:poly:def}{polynomial} identities by question 3's argument.

\textbf{5.} The $x_k$ are $n$ distinct simple roots and the
leading coefficient is $2^{n-1}$:
\[
T_n = 2^{n-1}\prod_{k=0}^{n-1}
\Bigl(X - \cos\frac{(2k+1)\pi}{2n}\Bigr) .
\]
Interlacing: the angles $0 < \frac{\pi}{2n} < \frac\pi n <
\frac{3\pi}{2n} < \frac{2\pi}n < \dots < \pi$ alternate between
the $y$-angles $\frac{k\pi}n$ and the $x$-angles
$\frac{(2k+1)\pi}{2n}$; since $\cos$ is strictly decreasing on
$\intcc0\pi$, the values interlace in the reverse order: $y_n <
x_{n-1} < y_{n-1} < \dots < x_0 < y_0$. Between two consecutive
extrema sits exactly one root, as a picture of $\cos n\theta$
suggests.

\textbf{6.} Induction with $2\cosh a\cosh b = \cosh(a + b) +
\cosh(a - b)$ (\cref{prop:b1:functions:hyprules}): $T_{n+1}(\cosh
t) = 2\cosh t\cosh nt - \cosh(n-1)t = \cosh(n+1)t$. For $x \geq
1$, write $x = \cosh t$ with $t \geq 0$; then $\eu^t = x +
\sqrt{x^2 - 1}$ and $\eu^{-t} = x - \sqrt{x^2 - 1}$, so
\[
T_n(x) = \cosh(nt)
= \frac{(x + \sqrt{x^2-1})^n + (x - \sqrt{x^2-1})^n}2 .
\]
For $x > 1$ the first term exceeds $\frac12(1)^n$ strictly and
grows geometrically: $T_n(x) > 1$.

\textbf{7.} De Moivre: $\cos n\theta = \Re\bigl((\cos\theta +
\iu\sin\theta)^n\bigr) = \sum_{2j \leq n}\binom n{2j}
\cos^{n-2j}\theta\,(\iu\sin\theta)^{2j}$, and $(\iu\sin\theta)^{2j}
= (-\sin^2\theta)^j = (\cos^2\theta - 1)^j$. Substituting $x =
\cos\theta$ and invoking question 3:
\[
T_n(x) = \sum_{0\leq 2j\leq n}\binom n{2j}x^{n-2j}(x^2 - 1)^j .
\]
For $n = 3$: $\binom30 x^3 + \binom32 x(x^2 - 1) = x^3 + 3x^3 -
3x = 4x^3 - 3x$, as in question 1.

\textbf{8.} $T_n(1) = \cos(n\cdot0) = 1$; $T_n(-1) = \cos(n\pi) =
(-1)^n$; $T_n(0) = \cos\frac{n\pi}2$, which is $0$ for odd $n$ and
$(-1)^{n/2}$ for even $n$.

\textbf{9.} Induction for $U_n(\cos\theta) =
\frac{\sin(n+1)\theta}{\sin\theta}$: true for $U_0 = 1$ and $U_1 =
2X$ ($\sin2\theta = 2\sin\theta\cos\theta$); the step is the
sum-to-product identity $\sin(n+2)\theta = 2\cos\theta\,
\sin(n+1)\theta - \sin n\theta$. Now differentiate
$T_n(\cos\theta) = \cos n\theta$ in $\theta$:
$-\sin\theta\,T_n'(\cos\theta) = -n\sin n\theta$, so for $\theta
\notin \pi\Z$:
\[
T_n'(\cos\theta) = n\,\frac{\sin n\theta}{\sin\theta}
= n\,U_{n-1}(\cos\theta) ,
\]
and the \omterm{def:b1:poly:def}{polynomials} $T_n'$ and $nU_{n-1}$, agreeing on
$\intoo{-1}1$, are equal.

\textbf{10.} $\abs{\sin(n+1)\theta} = \abs{\sin n\theta\cos\theta
+ \cos n\theta\sin\theta} \leq \abs{\sin n\theta} +
\abs{\sin\theta}$, and induction gives $\abs{\sin n\theta} \leq
n\abs{\sin\theta}$. Hence $\abs{U_{n-1}} \leq n$ on $\intoo{-1}1$
and $\abs{T_n'} = n\abs{U_{n-1}} \leq n^2$ there; at $\pm1$ the
bound extends by taking limits (or directly: $U_{n-1}(1) = n$ from
the recurrence, $U_n(1) = n + 1$ by induction, and parity gives
$U_{n-1}(-1) = (-1)^{n-1}n$). Thus $T_n'(1) = n^2$ and $T_n'(-1)
= (-1)^{n-1}n^2$: the bound $n^2$ is attained at the endpoints.

\textbf{11.} Differentiate $\sin\theta\,T_n'(\cos\theta) = n\sin
n\theta$ (question 9) with respect to $\theta$:
\[
\cos\theta\,T_n'(\cos\theta) - \sin^2\theta\,T_n''(\cos\theta)
= n^2\cos n\theta = n^2\,T_n(\cos\theta) .
\]
With $x = \cos\theta$ and $\sin^2\theta = 1 - x^2$: $x\,T_n' - (1
- x^2)T_n'' = n^2T_n$ on $\intcc{-1}1$, hence everywhere:
$(1 - x^2)y'' - xy' + n^2y = 0$ for $y = T_n$. Check for $T_2 =
2x^2 - 1$: $(1 - x^2)(4) - x(4x) + 4(2x^2 - 1) = 4 - 4x^2 - 4x^2
+ 8x^2 - 4 = 0$.

\textbf{12.} $\widetilde T_n$ is \omterm{def:b1:poly:def}{monic} (question 2) and
$\abs{\widetilde T_n} = 2^{1-n}\abs{T_n} \leq 2^{1-n}$ on
$\intcc{-1}1$, with $\widetilde T_n(y_k) = (-1)^k2^{1-n}$ at the
$n + 1$ points $y_k$ (\cref{exo:b1:poly:10}): the norm is exactly
$2^{1-n}$, attained with alternating signs.

\textbf{13.} $\widetilde T_n$ and $P$ are both \omterm{def:b1:poly:def}{monic} of degree
$n$, so the leading terms cancel: $\deg D \leq n - 1$. At $y_k$:
$D(y_k) = (-1)^k2^{1-n} - P(y_k)$, and $\abs{P(y_k)} \leq \norm
P_\infty < 2^{1-n}$ forces the sign of $D(y_k)$ to be that of
$(-1)^k2^{1-n}$, strictly.

\textbf{14.} $D$ changes sign between $y_{k+1}$ and $y_k$ for each
$k = 0, \dots, n-1$: by the intermediate value property, $D$ has a
root in each of these $n$ pairwise disjoint open intervals --- $n$
distinct roots for a nonzero \omterm{def:b1:poly:def}{polynomial} of degree $\leq n - 1$,
impossible. And $D = 0$ is impossible too (the norms differ).
Contradiction: no \omterm{def:b1:poly:def}{monic} $P$ of degree $n$ has $\norm P_\infty <
2^{1-n}$, which is Chebyshev's theorem.

\textbf{15.} Now $\abs{P(y_k)} \leq 2^{1-n}$ only, so $(-1)^k
D(y_k) = 2^{1-n} - (-1)^kP(y_k) \geq 2^{1-n} - \abs{P(y_k)} \geq
0$. Suppose $D(y_k) = 0$ at an interior $y_k$ ($0 < k < n$): then
$P(y_k) = (-1)^k2^{1-n}$, so $\abs P$ attains its supremum
$2^{1-n}$ at the interior point $y_k$, whence $P'(y_k) = 0$
(interior extremum); and $T_n'(y_k) = nU_{n-1}(y_k) = 0$ since
$\sin(n\cdot\frac{k\pi}n) = 0$ --- so $\widetilde T_n'(y_k) = 0$
too, and $D'(y_k) = 0$: $y_k$ is a root of $D$ of \omterm{def:b1:poly:derivative}{multiplicity} at
least $2$.

\textbf{16.} Count roots of $D$ with \omterm{def:b1:poly:derivative}{multiplicity}. Let $z$ be the
number of interior points $y_k$ with $D(y_k) = 0$ (each a double
root, by question 15) and $e \in \{0, 1, 2\}$ the number of
endpoints ($y_0$ or $y_n$) with $D = 0$ (each a simple root at
least). A gap $(y_{k+1}, y_k)$ whose two endpoints both have $D
\neq 0$ carries strictly alternating signs, hence an interior
root. Each vanishing interior point spoils at most its two
adjacent gaps, each vanishing endpoint at most one gap: at least
$n - 2z - e$ gaps still contribute one root each, all distinct
from the $y$-roots. Total: at least $(n - 2z - e) + 2z + e = n$
roots with \omterm{def:b1:poly:derivative}{multiplicity}, for a \omterm{def:b1:poly:def}{polynomial} of degree $\leq n - 1$:
so $D = 0$ and $P = \widetilde T_n$. The minimizer is unique.

\textbf{17.} The affine \omterm{def:b1:logic:map}{map} $t \mapsto x = \frac{a+b}2 +
\frac{b-a}2\,t$ is a bijection $\intcc{-1}1 \to \intcc ab$. If $P$
is \omterm{def:b1:poly:def}{monic} of degree $n$, then $Q(t) = P(x(t))$ is a \omterm{def:b1:poly:def}{polynomial} in
$t$ with leading coefficient $\bigl(\frac{b-a}2\bigr)^n$, and
$\sup_{\intcc ab}\abs P = \sup_{\intcc{-1}1}\abs Q$. The \omterm{def:b1:poly:def}{monic
polynomial} $Q/\bigl(\frac{b-a}2\bigr)^n$ has sup-norm $\geq
2^{1-n}$ (questions 13--14), so
\[
\sup_{\intcc ab}\abs P \geq \Bigl(\frac{b-a}2\Bigr)^n 2^{1-n}
= 2\Bigl(\frac{b-a}4\Bigr)^n ,
\]
with equality exactly for $P(x) = \bigl(\frac{b-a}2\bigr)^n
\widetilde T_n\bigl(t(x)\bigr)$ (question 16).

\textbf{18.} $\widetilde T_3 = \frac{T_3}4 = X^3 - \frac34X$;
$\widetilde T_3{}' = 3X^2 - \frac34$ vanishes at $\pm\frac12$.
Values: $\widetilde T_3(-1) = -\frac14$, $\widetilde
T_3(-\tfrac12) = \frac14$, $\widetilde T_3(\tfrac12) = -\frac14$,
$\widetilde T_3(1) = \frac14$: four alternating extrema of
absolute value $\frac14$ --- so $\norm{\widetilde T_3}_\infty =
\frac14$, and by Chebyshev's theorem no \omterm{def:b1:poly:def}{monic} cubic has smaller
sup-norm on $\intcc{-1}1$.

\textbf{19.} $\omega$ is \omterm{def:b1:poly:def}{monic} of degree $n + 1$, so
$\norm\omega_\infty \geq 2^{-n}$ by Chebyshev's theorem (degree
$n+1$), with equality if and only if $\omega = \widetilde T_{n+1}
= 2^{-n}T_{n+1}$ (question 16), i.e.\ if and only if the nodes
are the $n + 1$ roots of $T_{n+1}$. With Chebyshev nodes the error
factor $\norm\omega_\infty$ is $2^{-n}$ --- the least possible.

\textbf{20.} $5 \times 36^\circ = 180^\circ$, so $T_5(c) =
\cos180^\circ = -1$: $16c^5 - 20c^3 + 5c + 1 = 0$. Testing $x =
-1$: $-16 + 20 - 5 + 1 = 0$, and expanding confirms
\[
16x^5 - 20x^3 + 5x + 1 = (x + 1)\bigl(4x^2 - 2x - 1\bigr)^2 .
\]
Since $c = \cos36^\circ \neq -1$, $c$ is a root of $4x^2 - 2x -
1$, whose roots are $\frac{1 \pm \sqrt5}4$; as $c > 0$,
\[
\cos36^\circ = \frac{1 + \sqrt5}4 .
\]
Consistency: $\cos72^\circ = T_2(c) = 2c^2 - 1 = 2\cdot\frac{3 +
\sqrt5}8 - 1 = \frac{\sqrt5 - 1}4$, the value found in
\cref{exo:b1:complex:8}.

\textbf{21.} $\sqrt{1.1^2 - 1} = \sqrt{0.21} \approx 0.458$, so
$x + \sqrt{x^2-1} \approx 1.558$ and $(1.558)^{10} \approx 84.5$,
while $(1.1 - 0.458)^{10} \approx 0.01$: $T_{10}(1.1) \approx
\frac{84.5 + 0.01}2 \approx 42$. A \omterm{def:b1:poly:def}{polynomial} trapped in
$\intcc{-1}1$ on the interval has already grown past $40$ one
tenth beyond its edge: boundedness on a segment says nothing an
inch outside it.

\textbf{22.} In question 7's formula for $T_p$, the term $j = 0$
is $X^p$; every other term carries $\binom p{2j}$ with $0 < 2j <
p$ (note $2j \neq p$ since $p$ is odd), which is divisible by $p$
by the first step of the proof of \cref{thm:b1:arith:fermat}.
Hence every coefficient of $T_p - X^p$ is a multiple of $p$.
Checks: $T_3 - X^3 = 3X^3 - 3X = 3(X^3 - X)$; $T_5 - X^5 = 15X^5
- 20X^3 + 5X = 5(3X^5 - 4X^3 + X)$.

\textbf{23.} By question 17 with $\intcc ab = \intcc01$ and $n =
2$: minimal deviation $2\bigl(\frac14\bigr)^2 = \frac18$,
attained by $\bigl(\frac12\bigr)^2\widetilde T_2(2x - 1) =
\frac14\bigl((2x-1)^2 - \frac12\bigr) = x^2 - x + \frac18$. The
\omterm{def:b1:poly:def}{monic} quadratic closest to zero on $\intcc01$ is $x^2 - x +
\frac18$, with sup-norm $\frac18$.

\textbf{24.} (i) Rigidity --- a \omterm{def:b1:poly:def}{polynomial} with more roots than
its degree is zero --- powered the uniqueness principle
(question 3), the transfer of trigonometric identities to
\omterm{def:b1:poly:def}{polynomial} identities (questions 4, 7, 9, 11), and both
root-counting arguments of the extremality proof (questions 14,
16). (ii) The trigonometry of \cref{ch:b1:complex} (de Moivre,
sum-to-product) and the \omterm{def:b1:functions:hyperbolic}{hyperbolic functions} of
\cref{ch:b1:functions} supplied every identity behind the family;
the substitution $x = \cos\theta$ is the bridge. (iii) The
\omterm{def:b1:arith:divides}{divisibility} $p \mid \binom p{2j}$ from \cref{ch:b1:arith} turned
the coefficient formula into the \omterm{def:b1:arith:congruence}{congruence} of question 22.

\textbf{25.} Chebyshev's theorem converts an optimization over an
infinite-dimensional family (all \omterm{def:b1:poly:def}{monic polynomials}) into finite
combinatorics: a competitor better than $\widetilde T_n$ would
differ from it by a low-degree \omterm{def:b1:poly:def}{polynomial} forced to change sign
$n$ times --- one more root than its degree allows. The
equioscillation pattern is thus not a curiosity but the very
certificate of optimality, and the equality case sharpens
root-counting with multiplicities. The substitution $x =
\cos\theta$ deserves the last word: it transports the rigid,
discrete world of \omterm{def:b1:poly:def}{polynomials} into the periodic world of
trigonometry, where roots and extrema of $T_n$ are simply the
regular grid of $\cos n\theta$. The two analysis facts borrowed
--- the intermediate value property (question 14; proved in
\cref{ch:b1:continuity}) and the vanishing derivative at an
interior extremum (question 15; proved in \cref{ch:b1:derivative})
--- are exactly the tools those later chapters will hand back,
closing the loop.
\end{solution}
